
\documentclass[aspectratio=169]{beamer}

\usetheme{Madrid}
\usecolortheme{default}
\usefonttheme{professionalfonts}

\setbeamertemplate{navigation symbols}{}

\title[Combustion]{UCK 427 -- Combustion}
\subtitle{Combustion and Thermochemistry}
\author{Prof.\ Dr.\ Onur Tun\c{c}er}
\institute{Istanbul Technical University}
\date{Week 2}

\begin{document}

\begin{frame}
  \titlepage
\end{frame}

\begin{frame}{What Is Combustion?}

\begin{block}{Thermodynamic View}
Combustion is a chemical reaction process that converts 
chemical energy into thermal energy.
\end{block}

\vspace{0.3cm}

Energy release → temperature rise → pressure rise → work potential

\end{frame}

\begin{frame}{Energy Conversion Chain}

Chemical Bond Energy
\[
\Downarrow
\]
Internal Energy
\[
\Downarrow
\]
Thermal Energy
\[
\Downarrow
\]
Mechanical Work (e.g., turbine, rocket)

\end{frame}

\begin{frame}{Thermodynamic Description}

A reacting system state is defined by:

\[
(T, P, \{Y_i\})
\]

\pause

Composition is part of the thermodynamic state.
\end{frame}

\begin{frame}{Intensive vs Extensive Properties}

\textbf{Extensive:}
\[
U, H, S, V, m
\]

\textbf{Intensive:}
\[
T, P, \rho, u, h
\]

\pause

Scaling rule:

\[
X \rightarrow \lambda X
\quad
x = \frac{X}{m}
\]

\end{frame}

\begin{frame}{Why It Matters}

Reaction rates depend on:
\begin{itemize}
  \item Temperature
  \item Concentration
\end{itemize}

Total heat release depends on:
\begin{itemize}
  \item Total mass of fuel
\end{itemize}

\end{frame}

\begin{frame}{Ideal Gas Assumption}

For most combustion systems:

\[
p_i = \rho_i R_i T
\]

\pause

Total pressure:

\[
P = \sum_i p_i
\]

\end{frame}

\begin{frame}{Mixture Gas Constant}

\[
P = \rho R_{mix} T
\]

\[
R_{mix} = \sum_i Y_i R_i
\]

\pause

Note: $R_i = \frac{R_u}{M_i}$

\end{frame}

\begin{frame}{Internal Energy of Ideal Gas}

For ideal gases:

\[
u_i = u_i(T)
\]

\pause

\[
du_i = c_{v,i}(T)\, dT
\]

\end{frame}

\begin{frame}{Internal Energy of Mixture}

\[
U = \sum_i m_i u_i(T)
\]

\pause

Specific:

\[
u = \sum_i Y_i u_i(T)
\]

\end{frame}

\begin{frame}{Definition of Enthalpy}

\[
H = U + PV
\]

\pause

For ideal gases:

\[
h_i = u_i + R_i T
\]

\end{frame}

\begin{frame}{Mixture Enthalpy}

\[
h = \sum_i Y_i h_i(T)
\]

\pause

Differential form:

\[
dh = c_p^{mix} dT
\]

\[
c_p^{mix} = \sum_i Y_i c_{p,i}
\]

\end{frame}

\begin{frame}{Reference State}

Absolute enthalpy cannot be measured.

We define:

\[
h_i(T) = h_i(T_{ref}) + \int_{T_{ref}}^T c_{p,i}(T)\, dT
\]

Standard reference:

\[
T_{ref} = 298\,K
\]

\end{frame}

\begin{frame}{Enthalpy of Formation}

\begin{block}{Definition}
Energy change when 1 mole of compound forms from elements at standard state.
\end{block}

Example:

\[
C + O_2 \rightarrow CO_2
\]

\[
\Delta h_f^\circ = -393.5 \text{ kJ/mol}
\]

\end{frame}

\begin{frame}{Formation Enthalpy of Elements}

By convention:

\[
\Delta h_f^\circ(\text{O}_2) = 0
\]

\[
\Delta h_f^\circ(\text{N}_2) = 0
\]

\[
\Delta h_f^\circ(\text{H}_2) = 0
\]

\end{frame}

\begin{frame}{Total Species Enthalpy}

\[
h_i(T) = h_{f,i}^\circ + \int_{T^\circ}^{T} c_{p,i}(T)\, dT
\]

\pause

Two parts:
\begin{itemize}
  \item Chemical
  \item Sensible
\end{itemize}

\end{frame}

\begin{frame}{Heat of Reaction}

\[
\Delta h_r^\circ = 
\sum \nu_i h_{f,i}^\circ (products)
-
\sum \nu_i h_{f,i}^\circ (reactants)
\]

\pause

Exothermic:
\[
\Delta h_r^\circ < 0
\]

\end{frame}

\begin{frame}{Example: Methane Combustion}

\[
CH_4 + 2O_2 \rightarrow CO_2 + 2H_2O
\]

\pause

\[
\Delta h_r^\circ =
[-393.5 + 2(-241.8)]
-
[-74.8]
\]

\pause

\[
\approx -802 \text{ kJ/mol}
\]

\end{frame}

\begin{frame}{Physical Meaning}

Large negative $\Delta h_r^\circ$:

\begin{itemize}
  \item Strong chemical bonds formed
  \item Energy released
  \item Temperature increases
\end{itemize}

\end{frame}

\begin{frame}{Energy Conservation in Combustion}

Adiabatic system:

\[
H_{reactants} = H_{products}
\]

\pause

\[
\text{Chemical energy} \rightarrow \text{Sensible heating}
\]

\end{frame}

\begin{frame}{Why Flame Temperature Is Limited?}

Temperature rise limited by:

\begin{itemize}
  \item Heat capacity of the mixture
  \item Dissociation at high temperature
  \item Non-ideal effects
\end{itemize}

\end{frame}

\begin{frame}{Next Lecture}

\begin{itemize}
  \item Stoichiometry
  \item Equivalence ratio
  \item Adiabatic flame temperature derivation
  \item First computational example
\end{itemize}

\end{frame}


\end{document}
